\documentclass[11pt]{article}
\usepackage[margin=1in, headheight=14pt]{geometry}
\usepackage{amsmath, amssymb}
\usepackage{hyperref}
\usepackage{enumitem}
\usepackage{algorithm}
\usepackage{algpseudocode}
\usepackage{fancyhdr}
\usepackage{tcolorbox}
\usepackage{microtype}
\usepackage{tikz}
\usepackage{pgfplots}
\pgfplotsset{compat=1.18}
\usetikzlibrary{arrows.meta, positioning, calc, shapes.geometric}

\input{notation.tex}
\input{zk-notation.tex}
\input{environments.tex}
\input{lecture-macros.tex}

\pagestyle{fancy}
\fancyhf{}
\lhead{EECS 498/CSE 598: Zero-Knowledge Proofs}
\rhead{Winter 2026}
\lfoot{Lecture 23: Anonymous Credentials}
\rfoot{\thepage}
\renewcommand{\headrulewidth}{0.4pt}
\renewcommand{\footrulewidth}{0.4pt}

% Lecture-specific macros
\tcbuselibrary{skins, breakable}

\newtcolorbox{highlight}[1][]{
  colback=blue!5!white,
  colframe=blue!60!black,
  fonttitle=\bfseries,
  #1
}

\newtcolorbox{graybox}[1][]{
  colback=gray!8,
  colframe=gray!50,
  fonttitle=\bfseries,
  #1
}

\newcommand{\Issue}{\algo{Issue}}
\newcommand{\Present}{\algo{Present}}
\newcommand{\VerifyShow}{\algo{VerifyShow}}
\newcommand{\Show}{\algo{Show}}
\newcommand{\SigVerify}{\algo{VerifySig}}
\newcommand{\KVAC}{\ensuremath{\mathsf{KVAC}}}
\newcommand{\MAC}{\ensuremath{\mathsf{MAC}}}
\newcommand{\Idemix}{\ensuremath{\mathsf{Idemix}}}
\newcommand{\Zerocash}{\ensuremath{\mathsf{Zerocash}}}
\newcommand{\mDL}{\ensuremath{\mathsf{mDL}}}

\begin{document}

%-----------------------------------------------------------------------
% Title Block
%-----------------------------------------------------------------------
\begin{center}
  {\LARGE \textbf{EECS 498 / CSE 598: Zero-Knowledge Proofs}}\\[6pt]
  {\large Winter 2026}\\[4pt]
  {\large \textbf{Lecture 23: Anonymous Credentials}}\\[4pt]
  {\large Instructor: Paul Grubbs \quad
    \href{mailto:paulgrub@umich.edu}{\texttt{paulgrub@umich.edu}} \quad
    Beyster 4709}
\end{center}

\vspace{0.5em}
\hrule
\vspace{1em}

%-----------------------------------------------------------------------
% Agenda
%-----------------------------------------------------------------------
\section*{Agenda}
\begin{itemize}
  \item Malleability and stronger soundness notions
  \item Anonymous-credential functionality and security
  \item Classical constructions from signatures
  \item Keyed-verification anonymous credentials
  \item zkSNARK-based credentials, device binding, and revocation
\end{itemize}

\hrule
\vspace{1em}

%-----------------------------------------------------------------------
\section{Malleability and Stronger Soundness Notions}
%-----------------------------------------------------------------------

The simplified payment statements used to explain \Zerocash{} leave open an
important integrity issue.  If public transaction data such as an
\algo{info}-string or a public payee address is not cryptographically tied to
the proof statement, a miner can modify that public field without changing the
underlying zero-knowledge proof and thereby redirect the payment.

\begin{remark}
Ordinary knowledge soundness does not rule out this kind of transformation.  A
proof system can be knowledge sound while still allowing an adversary to turn a
valid proof into a different valid proof for a related statement.
\end{remark}

The \Zerocash{} design addresses this by authenticating the full transaction
blob.  One convenient way to express that idea is to generate a one-time
signature key for each transaction and verify a signature on the transaction
inside the proof system.

\begin{definition}[Simulation soundness]
A proof system is simulation sound if soundness continues to hold even in the
presence of valid proofs, including proofs generated by a simulator.
\end{definition}

\begin{remark}
Simulation extractability is a stronger knowledge-style variant of the same
theme.  These notions are strictly stronger than ordinary soundness or
knowledge soundness and are the relevant notions when proof malleability is a
real attack surface.
\end{remark}

%-----------------------------------------------------------------------
\section{Anonymous-Credential Functionality}
%-----------------------------------------------------------------------

An anonymous-credential system has three parties:
\begin{itemize}
  \item an \emph{issuer} that certifies user attributes,
  \item a \emph{user} that holds the credential,
  \item a \emph{relying party} or \emph{verifier} that checks a policy on the
        credential.
\end{itemize}

\begin{highlight}[title={Anonymous-Credential Interface}]
Let $m$ be the user attribute vector and let $\phi$ be a public validity
predicate on attributes.
\begin{enumerate}
  \item Issuance:
        \[
          C \leftarrow \Issue(sk_{\mathrm{iss}}, m).
        \]
  \item Presentation:
        \[
          \pi \leftarrow \Present(C, m, \phi).
        \]
  \item Verification:
        \[
          b \leftarrow \VerifyShow(pk_{\mathrm{iss}}, \phi, \pi) \in \{0,1\}.
        \]
\end{enumerate}
The verifier learns only that the credential is valid and that
\[
  \phi(m)=1.
\]
\end{highlight}

The natural example is age verification.  The issuer is the government, the
user holds attributes such as date of birth and license number, and the
relying party wants to evaluate a predicate such as ``age at least $21$.''  The
system should not reveal more than the predicate outcome.

\begin{definition}[Unlinkability]
An anonymous-credential system is unlinkable if the verifier cannot decide
whether two accepted presentation transcripts came from the same credential or
from two different credentials with attributes satisfying the same policy.
\end{definition}

Two security goals dominate:
\begin{itemize}
  \item \textbf{Unforgeability:} users cannot create valid credentials for
        unauthorized attributes.
  \item \textbf{Unlinkability:} repeated presentations cannot be linked back to
        one credential or one user.
\end{itemize}

Selective disclosure is the intermediate case in which the verifier learns a
chosen subvector of the attributes together with hidden proofs about the rest.

\begin{remark}
The issuer usually knows the user's actual attributes at issuance time.  The
privacy target is therefore primarily hiding from the relying party, including
the case in which the relying party already knows candidate attribute vectors.
\end{remark}

%-----------------------------------------------------------------------
\section{Historical Route to Modern Anonymous Credentials}
%-----------------------------------------------------------------------

The conceptual starting point is David Chaum's early work on blind signatures,
untraceable payments, and privacy-preserving authentication.  Those papers
identified the basic participants, the idea of unlinkable showing, and the
general tension between authentication and privacy long before efficient modern
zero-knowledge tools existed.

The first efficient anonymous-credential systems emerged from the
Camenisch--Lysyanskaya line of work.  Their central advance was not merely a
new credential abstraction, but a signature system equipped with efficient
protocols for obtaining a signature on hidden data and proving knowledge of a
signature without revealing the signed message.

Commercial systems such as \algo{Credentica} and \Idemix{} explored this
direction, but practical deployment lagged behind the ideas.  Later pairing-
based work further improved the efficiency of the underlying protocols,
including constructions associated with Boneh, Boyen, and Shacham.

%-----------------------------------------------------------------------
\section{Signatures with Efficient Protocols}
%-----------------------------------------------------------------------

The classical recipe for anonymous credentials starts from a signature scheme
supporting two extra protocol features:
\begin{enumerate}
  \item issuance of a signature on a hidden message,
  \item proof of knowledge of a signature on a message without revealing the
        signature itself.
\end{enumerate}

\begin{proposition}[Signature-to-credential template]
If a signature scheme supports hidden issuance and hidden proof of possession,
then it can be turned into an anonymous-credential system.
\end{proposition}

\begin{proof}[Proof sketch]
During issuance, the user obtains a signature on the attribute vector without
revealing the full vector to the issuer.  During presentation, the user proves
knowledge of a valid signature together with a proof that the disclosed
predicate holds on the signed attributes.  Hiding the signature itself prevents
the verifier from using it as a stable identifier across multiple
presentations.
\end{proof}

This observation explains why signatures and anonymous credentials are so
closely related.  A credential is a certified message, and a credential
presentation is a privacy-preserving proof of possession of that certified
message.

%-----------------------------------------------------------------------
\section{Keyed-Verification Anonymous Credentials}
%-----------------------------------------------------------------------

The \KVAC{} viewpoint replaces public verification by a \MAC-style relation
between the issuer and the verifier.  Let the shared key be
\[
  x = (x_0,x_1,\ldots,x_n) \in \F^{n+1},
\]
and let the credential message be
\[
  m = (m_1,\ldots,m_n).
\]

\begin{definition}[Algebraic \MAC]
Choose a random group element $u$.  The credential tag is
\[
  \sigma = (u,u'),
\]
where
\[
  u' = u^{x_0 + \sum_{i=1}^{n} x_i m_i}.
\]
\end{definition}

Public issuer parameters expose commitments to the secret-key coordinates and
support zero-knowledge proofs that a credential was formed correctly.  The
presentation protocols are then built from proofs of group-element relations
rather than from a general-purpose circuit.

At a high level, the issuance protocol computes the \MAC{} on the message and
provides a proof that the public issuer parameters, the shared key, and the
issued credential are mutually consistent.  The presentation protocol then
rerandomizes or blinds the credential and proves possession of a valid tag on
attributes satisfying the requested predicate.

\begin{remark}
The proof accompanying issuance prevents the issuer from handing out malformed
credentials that could later break unlinkability or selectively identify the
user.
\end{remark}

%-----------------------------------------------------------------------
\section{zkSNARK-Compatible Credentials}
%-----------------------------------------------------------------------

The classical constructions assume that the signature or \MAC{} system itself
has efficient hidden issuance and proof-of-possession protocols.  A different
approach is to start from a legacy signature scheme and prove the needed
relation with a zkSNARK.

Suppose a user already holds a structured message $m$ and a signature
\[
  \sigma
\]
under the issuer public key $pk_s$.

\begin{graybox}[title={Legacy-Compatible Proof-of-Possession Relation}]
For public key $pk_s$ and public policy $\phi$, the presentation proof shows
\[
  \exists\, m,\sigma
  \quad\text{such that}\quad
  \SigVerify(pk_s, m, \sigma) = 1
  \quad\text{and}\quad
  \phi(m)=1.
\]
\end{graybox}

This recipe is flexible because it can sit on top of an already deployed
credential format.  The privacy argument comes from the zero-knowledge proof
rather than from special-purpose algebraic show protocols.

The main cost is representation.  The proof system must arithmetize both
signature verification and the policy logic.  For modern credential formats,
the policy logic may include parsing a signed blob, extracting fields from a
\algo{JSON}-like structure, and evaluating predicates over those fields.

\begin{remark}
The efficiency requirement in the classical ``signatures with efficient
protocols'' recipe is substantial.  General-purpose signature schemes such as
\algo{ECDSA} are expensive inside ordinary \algo{R1CS} unless the construction
uses dedicated non-native arithmetic and parsing optimizations.
\end{remark}

This is the motivation for recent work on privacy-preserving age verification
from mobile driver's licenses.  The \mDL{} setting starts from a legacy signed
credential blob and uses a zkSNARK to prove possession of a valid signature and
the truth of an age predicate.

%-----------------------------------------------------------------------
\section{Device Binding and Revocation}
%-----------------------------------------------------------------------

A digital anonymous credential is easier to share than a physical identity
document unless the presentation capability is tied to a protected device.
Device binding addresses this by associating the credential with a secret key
stored in trusted hardware.  Credential presentation then requires the
device-resident secret key, so copying the outer credential blob is not enough.

\begin{definition}[Device binding]
Device binding is the property that valid credential presentations require a
secret key held in a protected hardware element associated with one device.
\end{definition}

The usual goal is not perfect personhood verification.  The goal is narrower:
make credential transfer materially harder than copying an unsigned data blob.

Revocation remains difficult because a revoked credential must stop working
without turning every presentation into a linkable identifier.  Revocation
lists and related techniques partially address this, but the tension between
privacy and revocation remains a central open design constraint in anonymous
credentials.

\bigskip
\hrule
\bigskip

\begin{thebibliography}{99}

\bibitem{Chaum82}
  D.~Chaum,
  ``Blind signatures for untraceable payments,''
  \emph{CRYPTO}, 1982.

\bibitem{CL04}
  J.~Camenisch and A.~Lysyanskaya,
  ``Signature schemes and anonymous credentials from bilinear maps,''
  \emph{CRYPTO}, 2004.

\bibitem{BBS04}
  D.~Boneh, X.~Boyen, and H.~Shacham,
  ``Short group signatures,''
  \emph{CRYPTO}, 2004.

\bibitem{CMZ14}
  M.~Chase, S.~Meiklejohn, and G.~Zaverucha,
  ``Algebraic MACs and keyed-verification anonymous credentials,''
  \emph{ACM CCS}, 2014.

\bibitem{FrigoShelat24}
  M.~Frigo and A.~Shelat,
  ``Anonymous credentials from ECDSA,''
  \emph{IACR ePrint 2024/2010}, 2024.

\end{thebibliography}

\end{document}
